% OLD ABSTRACT
%Refractory high entropy alloys are potential next–generation structural materials, yet the accurate prediction of phase stability in these systems is a significant challenge. 
%
%We present a study of two equimolar refractory high entropy alloys, NbMoTaTi--X (X = W, V). 
%
%\emph{Ab initio} calculations are performed to determine the vacancy exchange potential in both alloys. 
%
%The alloys are synthesized using the arc melting technique, followed by microstructural characterization.
%
%
%
%
%
%
%
%
% NEW ABSTRACT
\begin{comment}
Prediction of phase-dependent material transformations and phase transitions, such as embrittlement in cast austenitic stainless steels (CASS, remain a challenge. Here we experimentally study spinodal decomposition of \emph{bcc} delta ferrite in CASS and theoretically investigate \emph{bcc} Fe$_{1-x}$Cr$_{x}$ to simulate spinodal decomposition of \emph{bcc} delta ferrite into Cr-rich $\alpha$' and Fe-rich $\alpha$ phases. First principles \emph{ab initio} quantum elastodynamic (QED) simulations are performed to predict experimentally observed picosecond ultrasonic surface acoustic wave (SAW) peak splitting in alloys that have undergone spinodal decomposition. Emergence of new modes in the elastodynamically-unallowed, spinodal region for pristine allows serves as an absolute, non-destructive proof of stiffening towards embrittlement of cast austenitic stainless steels. Theoretical and experimental results indicate that SAW peak splitting predicts embrittlement in \emph{bcc} delta ferrite by spinodal decomposition, while elemental trends may also predict embrittlement in other alloys such as Cu-precipitation in reactor pressure vessel (RPV) steels. If these results hold true across other systems, SAW peak splitting can serve as a rapid predictor of the presence of a phase indicating embrittlement in the vast structural material compositional space. 
\end{comment}

Here we study spinodal decomposition of \emph{bcc} delta ferrite in cast austenitic stainless steels (CASS) by combining experimental observations with theoretical simulations of \emph{bcc} Fe$_{1-x}$Cr$_{x}$. Quantum elastodynamic (QED) simulations predicted ultrasonic surface acoustic wave (SAW) peak splitting, which serves as a non-destructive proof of stiffening towards embrittlement in CASS. Agreement of theoretical results of two observables, stiffness and thermal diffusivity, with experimental measurements confirm that SAW peak splitting could predict embrittlement in other alloys like Cu-precipitation in reactor pressure vessel (RPV) steels. Our findings imply that SAW peak splitting serves as a rapid predictor of embrittlement in a wide range of structural materials for assessment of their lifetimes in service.